\documentclass[12pt]{article}
\newcommand{\bottomMargin}{2cm}


\newcommand{\header}{
	ДИЗАЙН И АНАЛИЗ НА АЛГОРИТМИ - \\ ПРАКТИКУМ\\
	\vspace{0.1cm}
        Летен семестър, 2026 г., първо контролно\\
	\vspace{0.1cm}
}
\newcommand{\tl}{$0.3$ сек.}
\newcommand{\ml}{$256$ MB}

\input{structure.tex}

\begin{document}

\section{Задача К3. Експедиция}

В някои точки от дълъг прав път са разположени площадки за кацане. Известни са разстоянията от началото на пътя до всяка от площадките. Няколко експедиции искат да достигнат дадени точки от пътя, тръгвайки от началото му с хеликоптер. За целта, всяка експедиция трябва задължително веднъж да кацне на една от дадените площадки (не непременно най-близката до желаната точка), изразходвайки определено дадено количество гориво за хеликоптера (да отбележим, че ако например има само една площадка, намираща се точно в началото, хеликоптерът като излети, трябва да кацне на тази площадка).

След кацане на някоя от площадките, експедицията трябва да се придвижи пеша (напред или назад, ако желаната точка не съвпада с площадката) до желаната точка по пътя. При това придвижване за всяка единица от пътя се изразходва даден брой единици гориво.

Напишете програма \texttt{\textbf{expedition}}, която намира за всяка експедиция минималното количество гориво, необходимо за доставка на експедицията до посочената точка.

\subsection{Вход}
От първия ред на стандартния вход се въвеждат три цели числа $N$, $M$ и $C$ -- брой на площадките, брой на експедициите и разход на гориво за изминаване пеша на една единица дължина от пътя.

Следващите $N$ реда съдържат по две цели числа $A_i$ -- разстояние от началото на пътя до $i$-тата площадка и $B_i$ -- разходи на гориво за доставка на експедицията до $i$-тата площадка. Всички площадки са разположени в различни точки от пътя. 

От следващите $M$ реда се въвежда по едно цяло число $D_j$ -- разстоянието от началото на пътя до целта на $j$-тата експедиция.

\subsection{Изход}
За всеки полет на отделен ред на стандартния изход програмата трябва да изведе по едно цяло число -- минималният разход за превоз на експедицията.

\subsection{Ограничения}
\vspace{0.1em}
\begin{itemize}
	\item $1 \leq N\leq 10^5$
	\item $1 \leq M \leq 10^5$
	\item $1 \leq C \leq 10^9$
	\item $0 \leq A_i \leq 10^9$
	\item $0 \leq B_i \leq 10^9$
	\item $0 \leq D_j \leq 10^9$
\end{itemize}

\subsection{Оценяване:}
В около $33\%$ от тестовете $N, M \leq 1000$. 

\subsection{Пример}
\begin{table}[ht]
	\begin{tblr}{|X[15,l]|X[15,l]|X[70,l]|}
		\hline
		\textbf{Вход} & \textbf{Изход} & \textbf{Обяснение на примера} \\
		\hline
		{\texttt{3 2 1\\200 300\\300 100\\100 250\\150\\110}} & 
		{\texttt{250\\260}} &
		{Първата експедиция каца на втората площадка (на разстояние $300$ от началото на пътя), изразходвайки $100$ единици гориво и след това преминава до точка $150$, изразходвайки още $150$ единици гориво.\\[1em] Втората експедиция каца на третата площадка (на разстояние $100$ от началото на пътя), изразходвайки $250$ единици гориво и след това отива до точка $110$, като изразходва още $10$ единици гориво.} \\
		\hline
	\end{tblr}
\end{table}

\end{document}